% Folder 136, batch 07 — archivists' pages 121–140.
% Pass: Opus 5.5 (claude-opus-5-5), 2026-10-03 — first pass, unchecked against the pages by a human.
%
% Pages carrying his hand on the divided-power De Rham material: 122, 124,
% 126, 128, 130, 132, 134, 136, 138, 140. Page 121 is a blank coloured sheet;
% the odd pages 123–139 are printed proofs of an unrelated SGA 7 exposé
% (versos), skipped without description.
\documentclass[11pt,a4paper]{article}
\input{../preamble/grothendieck}

\folder{136}
\batch{7}
\pages{121}{140}
\foldertitle{Complexe de De Rham à puissance divisée [conférence de 1976 à l’IHÉS] : notes manuscrites (s.d.).}
\dating{[à partir de 1975-1976]}
\watermark{Édition de démonstration}

\begin{document}

\section{Complexes \uncertain{duaux} de De Rham à p.\,d.}

\page{122}
\note{titre encadré en haut à droite de la page, de sa main ; le mot souligné
est d'une lecture incertaine. La page ouvre une nouvelle rédaction.}

Soient

$k$ anneau \uncertain{commutatif} \add{ass. unitaire comm.}

\struck{\ill{}} $s$ une $k$-coalgèbre \add{$[s\to s\otimes_k s]$}
\add{ass. unitaire comm.} [donc \add{$V_k(s)$ :} $k'\mapsto s^{k'}=\operatorname{Hom}_k(s,k')$
est un foncteur à valeurs dans les anneaux \uncertain{commut.} \add{ass. un. comm.}]

\struck{\ill{}}

$\bar s$ une sous-$k$-coalgèbre \add{ass. unit. comm.} [donc
\add{$V_k(\bar s)$ :} $k'\mapsto \bar s^{k'}=\operatorname{Hom}_k(\bar s,k)$ est une
$\mathcal{O}_k$-alg. ass. un. comm. et $s^{k'}\to\bar s^{k'}$ est un hom.
de $k'$-alg. ass. unit. comm., i.e. $V_k(s)\to V_k(\bar s)$ hom. de
$\mathcal{O}_k$-alg.]
\note{une double flèche verticale relie $\bar s$ et $s^+$ dans la colonne de
gauche.}

$s^+=s/\bar s$ le $k$-module quotient [i.e. une coalgèbre ass. comm. mais
\ill{}] \quad $[\,s^{+k'}\simeq\operatorname{Ker}(s^{k'}\to\bar s^{k'})\subset s^{k'}$ est un
idéal de $s^{k'}$, donc $V_k(s^+)\overset{\mathrm{df}}{=}V_k(s)^+$ est un Idéal de
$V_k(s)\,]$

$\gamma_*$ un système de copuissances divisées sur $(s,s^+)$
[donc on a sur $V_k(s^+)=V_k(s)^+$ un syst. de puiss. divisées
\[
  \gamma_i : V_k(s^+)\longrightarrow V_k(s^+) \qquad i\geqslant 1
\]
défini par des homom. de $k$-mod.
\[
  \gamma_i^{*} : \operatorname{Sym}_k(s^+)\longrightarrow \operatorname{Sym}^i_k(s^+)
\]
\struck{\ill{}} i.e. homom. \uncertain{$k$-lin.}
\[
  \gamma_i^{*} : s^+\longrightarrow \operatorname{Sym}^i_k(s^+)
\]
(qui, au vu des axiomes, sont en fait)
\[
  \boxed{\gamma_i^{*} : s^+\longrightarrow \operatorname{Sym}^i_k(s^+)\quad (i\geqslant 1)}
\]
avec des axiomes divers

$t : s^+\to k$ un $k$-hom., i.e.
$t\in V_k(s^+)(k)=V_k(s)^+(k)$.

Alors on définit un objet simplicial dans la cat. des foncteurs
\[
  (k\text{-alg. comm.})\longrightarrow
\]
(algèbres graduées alternées différentielles (diff. de degré $1$) à
\add{avec idéaux} \struck{puiss. divisées})

$k'\mapsto$ \ill{} des \struck{foncteurs} schémas affines en algèbres graduées
alternées différentielles \add{homogènes} (diff. de degré $+1$) dans
\ill{}, à puissances divisées (\ill{} \ill{}), les puissances divisées étant
compatibles aux graduations et à la différentielle

\page{124}
\[
  \mathcal{A}^{*}_{*}(s,s^+,\gamma_*,t)(k')
  =\mathcal{A}^{*}_{*}\bigl(s^{k'},(s^{k'})^+,\gamma_*^{k'},t^{k'}\bigr)
\]
\note{sous $s^{k'}$ et $(s^{k'})^+$ il écrit respectivement
$V_k(s)(k')$ et $V_k(s^+)(k')$ ; sous le membre de gauche, en accolade :}
schéma affine en alg. graduées (alternées différentielles (dif. de
degré $+1$), avec idéal hom. à p.\,d. (comp. au reste))

On a
\[
  \mathcal{A}_*(s,s^+,\gamma_*,t)(k')(I)\simeq
  \widehat{\Gamma}_{s^{k'}_*}\{\{k'^{I}\}\}\otimes_{k'}
  \Lambda^{*}(k'^{I}/k')\big/\text{idéal dif. engendré par }\textstyle\sum X_i-t
\]

\[
  \mathcal{A}_n(s,s^+,\gamma_*,t)(k')\simeq
  s^{k'}\{\{X_0,\dots,X_n\}\}[dX_0,\dots,dX_n]\big/
  \bigl(\textstyle\sum X_i-t,\ \sum dX_i\bigr)_{\mathrm{pd}}
\]
\begin{align*}
  &\simeq s^{k'}\{\{X_0,\dots,X_{n-1}\}\}[dX_0,\dots,dX_{n-1}]\\
  &\simeq \prod\bigl(s^{k'}\bigr)^{\mathbb{N}^n\times\mathfrak{P}([0,n-1])}
  \qquad\text{comme foncteur en } k'\\
  &\simeq \operatorname{Hom}_k(s,k')^{\mathbb{N}^n\times\mathfrak{P}([0,\dots])}
  \qquad\text{comme foncteur en } k'\\
  &\simeq \operatorname{Hom}_k\bigl(s^{(\mathbb{N}^n\times\mathfrak{P}([0,\dots]))},k'\bigr)\\
  &= V\bigl(s^{(\mathbb{N}^n\times\mathfrak{P}([0,\dots]))}\bigr)(k')\\
  &= V\bigl(s\otimes_k k[Y_0,\dots,Y_{n-1}]\otimes_k
     {\textstyle\Lambda^{*}}(y_0,\dots,y_{n-1})\bigr)(k')
\end{align*}
\note{une flèche part de l'exposant $\mathbb{N}^n\times\mathfrak{P}([0,n-1])$ vers
« avec des \ill{} » ; le $\mathfrak{P}$ (parties) est écrit barré d'un trait,
comme il le fait d'ordinaire.}

\[
  \mathcal{A}_n(s,s^+,\gamma_*,t)(k')\big/\mathcal{A}_n(\ )(k')^+
  \simeq s^{k'}\big/(s^{k'})^+\hookrightarrow \bar s^{k'}
\]
\[
  \mathcal{A}_n(\ )(k')\longrightarrow \bar s^{k'}=V(\bar s)(k')
\]
\[
  V\bigl(s\otimes_k k[Y_0,\dots,Y_{n-1}]\otimes_k{\textstyle\Lambda^{*}}(y_0,\dots,y_{n-1})\bigr)(k')
  \longrightarrow V(\bar s)(k')
\]
\[
  s\otimes_k k[Y_0,\dots,Y_{n-1}]\otimes_k{\textstyle\Lambda^{*}}(y_0,\dots,y_{n-1})
  \longleftarrow \bar s
\]
\note{à droite, $\bar s$ est relié par une flèche à $s$, et une longue flèche
revient de $s$ vers le produit tensoriel : la flèche
$\bar s\to s\otimes_k\cdots$ se factorise par $\bar s\to s$.}

On voit donc, \ill{} \ill{}
\[
  \mathcal{A}^{n}_{*}(s,s^+,\gamma_*,t)
  =s\otimes_k k[Y_0,\dots,Y_{n-1}]\otimes_k{\textstyle\Lambda^{*}}(y_0,\dots,y_{n-1})
\]

\page{126}
a) Chacun des $\mathcal{A}^n_*(-)$ est muni d'une structure de $k$-coalgèbre
\add{(alternée)} graduée codifférentielle (codifférentielle de degré $-1$),
avec une immersion $\bar s\hookrightarrow\mathcal{A}^n_*(-)$ qui est compatible avec
tout (provenant de $\bar s$, à la graduation nulle, et à la différentielle
nulle), et \struck{une} syst. de copuissances divisées, sur
$\mathcal{A}^n(\ )^+=\mathcal{A}^n(\ )/\bar s$, compatibles aux graduations et
différentielle.

b) Pour $n$ variable, les $\mathcal{A}^n(\ )$ forment un objet
$\tfrac12$\add{co}simplicial (et $=$ simplicial) $\mathcal{A}^*(s,s^+,\gamma_*,t)$
dans la cat. des (coalg. \ill{} …) \add{objets, munis des structures}
\uncertain{dites} dans a).

c) On reconstitue $\mathcal{A}_*(s,s^+,\gamma_*,t)$ \add{avec ses} structures à
l'aide de $\mathcal{A}^*(s,s^+,\gamma_*,t)$, comme
\[
  \mathcal{A}_*(s,s^+,\gamma_*,t)=V_k\bigl(\mathcal{A}^*(s,s^+,\gamma_*,t)\bigr)
\]
\add{$=k'\mapsto\mathcal{A}_*(s,s^+,\gamma_*,t)(k')$}.

d) Soit $\mathcal{X}$ un ens. $\tfrac12$simplicial, considérons
\[
  k'\longmapsto\operatorname{Hom}\bigl(\mathcal{X},\mathcal{A}_*(-)(k')\bigr)
  \overset{\mathrm{df}}{=}
  C^{*}_{DR}\bigl(\mathcal{X};s^{k'},(s^{k'})^+,\gamma_*,t\bigr)
\]
\begin{align*}
  &\simeq \varprojlim_{\sigma\text{ dans }\Delta/\mathcal{X}}
     \operatorname{Hom}\bigl(\sigma,\mathcal{A}_*(\ )(k')\bigr)\\
  &\simeq \varprojlim_{\sigma\text{ dans }\Delta/\mathcal{X}}
     \mathcal{A}_{\dim\sigma}(\ )(k')
     \qquad\bigl(=\operatorname{Hom}_k(\mathcal{A}^{\dim\sigma}_*(-),k')\bigr)\\
  &\simeq \operatorname{Hom}_k\Bigl(\varinjlim_{\sigma\text{ dans }\Delta/\mathcal{X}}
     \mathcal{A}^{\dim\sigma}(-),k'\Bigr)\\
  &\simeq \operatorname{Hom}_k\bigl(C^{DR}_*(\mathcal{X};s,s^+,\gamma_*,t),k'\bigr)
\end{align*}
\note{en face de $\Delta/\mathcal{X}$ il précise : « cat. des simplexes types
$(\Delta_i)_{i\geqslant0}$, $\Delta_i=[0,i]$ ». Aux deux dernières lignes
l'indice d'origine de $\mathcal{A}$ est biffé et remplacé par
$\dim\sigma$ en exposant.}

\page{128}
\[
  \boxed{C^{DR}_*(\mathcal{X};s,s^+,\gamma_*,t)\overset{\mathrm{df}}{=}
  \varinjlim_{\sigma\text{ dans }\Delta/\mathcal{X}}
  \mathcal{A}^{\dim\sigma}(s,s^+,\gamma_*,t)}
\]

Ici $C^{DR}_*(\mathcal{X};s,s^+,\gamma_*,t)$ est une $k$-coalgèbre graduée
(alternée) codiff. (codiff. de degré $-1$), \struck{à copuissances}
\add{avec une structure} homom. $\bar s\to C^{DR}_*(\mathcal{X},-)$ (pas
nécessairement injectif : si $\mathcal{X}=\varnothing$,
$C^{DR}_*(\mathcal{X},-)=0$ …) et des copuiss. divisées sur
$C^{DR}_*(\mathcal{X},-)_+=$ \ill{}
$(\bar s\to C^{DR}_*(\mathcal{X},-))$. Il en \uncertain{résulte des}
structures, sur $C^*_{DR}(\mathcal{X},s^{k'},\dots)$ \uncertain{contravariantes},
\ill{} varie sur $k$, \struck{\ill{} varie avec $\mathcal{X}$}.

\marginal{écrit en oblique dans la marge gauche :
$\mathcal{X}\mapsto C^{DR}_*(\mathcal{X};\dots)$ \uncertain{dépend de}
$\mathcal{X}$ \uncertain{en commute} aux $\varinjlim$ quelconques
(\uncertain{vu sa défin.} \ill{})}

e) Si $(s,s^+,\gamma_*,t)$ varie sur $k$,
$(s^{k'},s^{+k'},\gamma_*^{k'},t^{k'})$ en dépend de façon contravariante,
donc
\[
  C^*_{DR}\bigl(\mathcal{X};s^{k'},s^{+k'},\gamma_*,t^{k'}\bigr)
  =\operatorname{Hom}\bigl(C^{DR}_*(\mathcal{X},-),k'\bigr)
\]
aussi, et $C^{DR}_*(\mathcal{X};s,s^+,\gamma_*,t)$ dépend de façon
covariante de $(s,s^+,\gamma_*,t)$, ce qui est aussi évident sur l'écriture
plus haut $C^{DR}_*(\mathcal{X},-)=\varinjlim\dots$, compte tenu que
$\mathcal{A}^*(s,s^+,\gamma_*,t)$ est covariant en $(s,s^+,\gamma_*,t)$.
De plus, \struck{il} comme fonction en $(s,s^+,\gamma_*,t)$, il commute aux
$\varinjlim$ filtrantes. \add{Commutation à l'extension de l'anneau de base.}

f) Supposons que la coalgèbre $s$ soit graduée (structure de coalgèbre et
graduation compatibles), que $\bar s$, sous-alg. \uncertain{itou},
$\bar s\to s$ compatible avec les gradu., les $\gamma_*$ \uncertain{préd.} sur
$s^+$ compatibles \add{avec} les graduations et $t : s^+\to k$ \add{homom.}
de degré $\delta$. Alors \struck{les} \add{$\forall k'$}
$\bigl(s^{k'},(s^{k'})^+,\gamma_*^{k'},t^{k'}\bigr)$

\page{130}
est une \add{alg.} \struck{à p. div.} graduée, avec idéal $(s^{k'})^+$
homog. à p. div. compatibles aux graduations, et $t^{k'}$ est un élément de
degré $\delta$ \uncertain{dedans}.

Alors les
\[
  \mathcal{A}_n(\ )(k')=s^{k'}\{\{X_0,\dots,X_n\}\}[dX_0,\dots,dX_n]
  \big/\bigl(\textstyle\sum X_i-t,\ \sum dX_i\bigr)_{\mathrm{pd}}
\]
est \add{une alg.} bigraduée, \add{(\uncertain{induite par} $s^{k'}$)}
(2ème degré) et le degré total défini par
\[
  \left\{
  \begin{aligned}
    &\deg\mathrm{tot}\,X_r=\deg\mathrm{tot}\,dX_r=\delta
      \quad\text{\struck{degré}}\\
    &\deg\mathrm{tot}\,X_r^{(n)}=\delta,\ \ \deg\mathrm{tot}_{\mathcal{A}_n}\lambda
      =\deg\mathrm{tot}_{s^{k'}}\lambda \qquad \lambda\in s^{k'}
  \end{aligned}\right.
\]
La diff. \struck{\ill{}} respecte le degré total, \ill{} une \ill{} graduée
\ill{} \uncertain{comme on l'a dit} ; les \uncertain{puiss. divisées} sont
compatibles aux deux degrés.

Cette structure s'\uncertain{explicite} \struck{\ill{}}
fonctorielle en $k'$ ; donc \ill{} (deg tot), et dans \ill{} :
$\mathcal{A}^*(s,s^+,\gamma^*,t)$.

Elle passe \ill{} \ill{} ainsi aux $C^{DR}_*(\mathcal{X};s,s^+,\gamma^*,t)$, et
on retrouve la structure \ill{} sur
$C^*_{DR}\bigl(\mathcal{X};s^{k'},(s^{k'})^+,\gamma^{*k'},t^{k'}\bigr)$ en
prenant le $\operatorname{Hom}_k$ \ill{} de
$C^{DR}_*(\mathcal{X};s,s^+,\gamma^*,t)$ dans $k'$.

g) On prend maintenant, pour
\begin{align*}
  s&=s_k=k[u] &&\text{donc } s^{k'}=k'\{\{T\}\}\\
  \bar s&=\bar s_k=k &&\text{donc } \bar s^{k'}=k'
\end{align*}
avec structure de coalg. évidente, \add{donnant sur $s^{k'}$ la structure
connue},
\[
  \bar s\longrightarrow s \ \text{l'inclusion}\ k\to k[u]
  \qquad\text{donc } s^{k'}\to\bar s^{k'}\ \text{l'augm. ordinaire}\
  k'\{\{T\}\}\to k'
\]
\[
  s^+=k[u]/k\simeq k[u]^+
\]

\page{132}
On prend syst. de copuiss. divisées sur $k[u]^+=s^+$ \add{comme}
correspondant à celles \struck{\ill{}} sur $k'\{\{T\}\}^+$.

On prend $t\in k[\{T\}]$ donné par $t=T$, i.e.
\[
  t(u^n)=\begin{cases}1 & \text{si } n=1\\ 0 & \text{si } n\neq1\end{cases}
\]

On a donc ici
\begin{align*}
  \mathcal{A}_*(s,s^+,\gamma^*,t)(k')
  &\simeq\mathcal{A}_*\bigl(s^{k'},s^{+k'},\gamma^{*k'},t^{k'}\bigr)
   \simeq\mathcal{A}_*\bigl(k'\{\{T\}\},k'\{\{T\}\}^+,\gamma^{k'},T\bigr)\\
  &\simeq\mathcal{A}_*(k')^{\wedge}
\end{align*}

On pose simplement
\[
  \text{\struck{$\mathcal{A}_k^{*}$}}\ \mathcal{A}^{*}_{k*}
  =\mathcal{A}^{*}_{*}(s_k,\bar s_k,\gamma_k^{*},t_k)
\]
de sorte que
\[
  \mathcal{A}_*(k')^{\wedge}\simeq\operatorname{Hom}_k(\mathcal{A}^{*}_{k},k')
\]
avec toutes les structures. Notons qu'ici on a $s$, $\bar s$ et
$\bar s\to s$ \add{gradués} (\ill{} \ill{} de $k[u]$), et dans le cas \ill{}
les graduations \ill{} sur les $k'\{\{T\}\}$), et $t=T$ de degré $1$, donc
\struck{on a} $\mathcal{A}^{*}_{k}$ admet une bigraduation. \add{Cela}
\struck{$\mathcal{A}^*_k$} donne lieu aussi à
\[
  \text{\struck{$C^{DR}_*(\mathcal{X};k)$}}\quad
  C^{DR,\mathrm{pd}}_*(\mathcal{X},k)\overset{\mathrm{df}}{\simeq}
  C^{DR}_*\bigl(\mathcal{X},k[u],k[u]^+,\gamma^{*}_k,t_k\bigr)
\]
(avec les \uncertain{autres structures}) et on retrouve par
$\operatorname{Hom}_k\bigl(C^{DR,\mathrm{pd}}_*(\mathcal{X},k),k'\bigr)$ les
structures connues (y compris \add{la} bigraduation) sur le complexe de
cochaînes de De Rham $C^*_{DR,\mathrm{pd}}(\mathcal{X},k')$.

En fait tous les $C^{DR,\mathrm{pd}}_*(\mathcal{X},k)$ \uncertain{se
déduisent}

\page{134}
\uncertain{se} déduisent de
\[
  C^{DR,\mathrm{pd}}_*(\mathcal{X})\overset{\mathrm{df}}{=}
  C^{DR,\mathrm{pd}}_*(\mathcal{X};\mathbb{Z})
\]
par la formule
\[
  C^{DR,\mathrm{pd}}_*(\mathcal{X},k)\simeq
  C^{DR,\mathrm{pd}}_*(\mathcal{X})\otimes_{\mathbb{Z}}k ,
\]
\uncertain{d'où} \ill{}
\[
  C^{*}_{DR,\mathrm{pd}}(\mathcal{X},k)\simeq
  \operatorname{Hom}\bigl(C^{DR,\mathrm{pd}}_*(\mathcal{X}),k\bigr)
\]

On définit de même, pour tout groupe abélien $M$ \uncertain{et
fonctoriellement en} $M$
\begin{align*}
  C^{DR,\mathrm{pd}}_*(\mathcal{X};M)&\overset{\mathrm{df}}{=}
  C^{DR,\mathrm{pd}}_*(\mathcal{X})\otimes_{\mathbb{Z}}M\\
  C^{*}_{DR,\mathrm{pd}}(\mathcal{X};M)&\overset{\mathrm{df}}{=}
  \operatorname{Hom}\bigl(C^{DR,\mathrm{pd}}_*(\mathcal{X}),M\bigr)
\end{align*}
et les structures sur $C^{DR,\mathrm{pd}}_*(\mathcal{X})$ \uncertain{impliquent}
\uncertain{une partie} des structures multiplicatives pour ces foncteurs en
$M$ …

\page{136}
\note{le haut de la page, jusqu'à « Dualisation… », est biffé de longs traits
obliques ; il ne prolonge pas visiblement la p. 134 et semble appartenir à
un autre fil (faisceaux sur un ensemble simplicial).}
\struck{ce qui se voit directement, \ill{}}
\[
  \text{\struck{$\tau^{\alpha}(F)=|\alpha|^{*}(\sigma^{\alpha}(F))\ \dots$}}
\]
\struck{Comme les $\sigma^{\alpha}$ commutent aux $\varinjlim$ finies et aux
$\varprojlim$ quelconques}
\[
  \text{\struck{$\mathcal{X}'=\varprojlim\mathcal{X}_{\mathcal{U}}\to X,\quad
  \mathcal{X}_{\mathcal{U}}\to\mathcal{X},\quad
  \varinjlim F/X_i,\quad \Delta/\mathcal{X}_i,\ \Delta/\mathcal{X}_j,\quad
  \mathcal{X}_i\hookrightarrow\mathcal{X}_j,\quad
  \mathcal{F}_i\leftarrow\mathcal{F}_j$}}
\]
\note{ces formules sont disposées en schéma sous les traits de biffure ; on
n'en donne que les termes.}

\subsection{Dualisation des homom. de Stokes}

\[
  \pi_n : \mathcal{A}^n_n(s,s^+,\gamma^*,t)(k')
  =\mathcal{A}^{\wedge}_n\bigl(s^{k'},(s^{k'})^+,\gamma^{k'},t^{k'}\bigr)^n
  \longrightarrow s^{k'}
\]
\note{sous le membre de gauche, relié par $=$ :
$V\bigl(s[Y_0,\dots,Y_n][y_0,\dots,y_n]^{n}\bigr)(k')$, le crochet extérieur
surchargé ; le but $s^{k'}$ est glosé $V(s)(k')$. En face de la flèche
verticale : « $\pi_n$ \uncertain{s-homom.} ».}

\[
  (\mathcal{A}^{*}_{*})_n\qquad
  s\xrightarrow{\ \kappa_n\ } s[Y_0,\dots,Y_n][y_0,\dots,y_n]^{(n)}
\]

1) c'est l'identité si $n=0$

2) Stokes : \uncertain{commutativité} dans

\begin{tikzcd}
  s \arrow[r, "\kappa_n"] \arrow[d, "\pi_{n-1}"'] & (\mathcal{A}^{*})_n \arrow[d, "d_{\mathcal{A}^*}"] \\
  (\mathcal{A}^{*})_{n-1} \arrow[r, "\partial_{ss}"'] & (\mathcal{A}^{*})_{n-1}
\end{tikzcd}

\note{dans le diagramme, la flèche verticale de gauche est notée
$\pi_{n-1}$ (lecture incertaine : $\kappa_{n-1}$ attendu), et la source en bas
à gauche porte un exposant surchargé.}

\[
  K^n_{\ill{}}(d\omega)=K^{n-1}_{\ill{}}(\partial\omega)
  \quad\Longleftrightarrow\quad
  d_{\mathcal{A}^*}\,\kappa_n=\partial_{ss}\,\kappa_{n-1}
\]
\note{sous $d_{\mathcal{A}^*}$ : « degré cot. $+1$ » ; sous $\partial_{ss}$ :
« deg ss $+1$ ». À gauche, un premier essai de la même formule est biffé :
\struck{$d\kappa_n=\partial\kappa_{n+1}\partial_{ss}$, degré $+1$}.}

\struck{$s\otimes k'$}
\[
  C^{\mathrm{sing}}_*(\mathcal{X},s)\longrightarrow
  C^{DR}_*(\mathcal{X};s,s^+,\gamma_*,t)
\]
\[
  C^{\mathrm{sing}}_*(\mathcal{X},k[u])\longrightarrow
  C^{DR,\mathrm{pd}}_*(\mathcal{X},k)
\]

\page{138}
\[
  \mathcal{A}^{0}_n\simeq k\{\{T\}\}\{\{X_0,\dots,X_n\}\}\big/
  \bigl(\textstyle\sum X_i-T\bigr)_{\mathrm{pd}}
  \xleftarrow{\ \text{épi}\ } k\{\{T\}\}\{\{X_0,\dots,X_n\}\}
\]
\[
  \simeq k\{\{X_0,\dots,X_n\}\}
  \longleftarrow k\{\{T,X_0,\dots,X_n\}\}
  \qquad X_i\mapsto X_i,\ \ \textstyle\sum_{0\leqslant i\leqslant n}X_i\leftarrow T
\]
\note{l'exposant $0$ de $\mathcal{A}$ est glosé « deg cot ».}

dualisant

\[
  (\mathcal{A}^n)_0\simeq k[U+Y_0,U+Y_1,\dots,U+Y_n]
  \xrightarrow{\ \text{mono}\ } k[U][Y_0,\dots,Y_n]\simeq k[U,Y_0,\dots,Y_n]
\]
\[
  k[Z_0,Z_1,\dots,Z_n]\longrightarrow k[U,Y_0,\dots,Y_n],\qquad
  Z_i\mapsto U+Y_i
\]
[l'image est l'intersection des noyaux des op. diff.
$\bigl(\tfrac{\partial}{\partial Y_0}+\dots+\tfrac{\partial}{\partial Y_n}
-\tfrac{\partial}{\partial U}\bigr)^{(i)}$]

\[
  \bigl\langle X_0^{(i_0)}\cdots X_n^{(i_n)},\,Z_0^{j_0}\cdots Z_n^{j_n}\bigr\rangle
  =\prod_{0\leqslant\alpha\leqslant n}\delta_{i_\alpha,j_\alpha}
  =\begin{cases}0 & \text{si } (i_\alpha)\neq(j_\alpha)\\
  1 & \text{si } (i_\alpha)=(j_\alpha)\end{cases}
\]
\[
  \bigl\langle T^{(\lambda)}X_0^{(i_0)}\cdots X_n^{(i_n)},\,
  U^{\mu}Y_0^{j_0}\cdots Y_n^{j_n}\bigr\rangle
  =\delta_{\lambda,\mu}\prod_\alpha\delta_{i_\alpha,j_\alpha}
  =\begin{cases}0 & \text{si } (\lambda,(i_\alpha))\neq(\mu,(j_\alpha))\\
  1 & \text{si } \dots=\dots\end{cases}
\]

\[
  \mathcal{A}_n^{*=*}\simeq
  {\textstyle\Lambda}^{*}_k[dX_0,\dots,dX_n]\big/\bigl(\textstyle\sum dX_i\bigr)
  \xleftarrow{\ \text{épi}\ }{\textstyle\Lambda}^{*}_k[dX_0,\dots,dX_n]
\]
\note{l'indice $*=*$ est glosé « degré tot. $=$ deg cot, i.e. deg compl. $=0$ ».}

dualisant
\[
  (\mathcal{A}^n)_{*=*}\simeq
  \bigl({\textstyle\Lambda}^{*}_k[y_0,\dots,y_n]\bigr)^{!}
  \simeq{\textstyle\Lambda}^{*}_k\bigl([y_0,\dots,y_n]^{!}\bigr)
  \longrightarrow{\textstyle\Lambda}^{*}_k[y_0,\dots,y_n]
\]
où $[y_0,\dots,y_n]^{!}$ est le ss-esp. de $ky_0\oplus\dots\oplus ky_n$ formé
des vect. dont la somme des coord. $=0$ ; l'image est le noyau de $i_X$, où
$X$ est la forme sur $ky_0\oplus\dots\oplus ky_n$ « somme des coordonnées » et
où $i_X$ est le produit int\supplied{érieur}.

\emph{NB} On a un iso can.
\[
  {\textstyle\Lambda}^{n}_k\bigl([y_0,\dots,y_n]'\bigr)\simeq
  {\textstyle\Lambda}^{n+1}_k\bigl([y_0,\dots,y_n]\bigr)\simeq k
\]
déduit de la suite exacte
\[
  0\longrightarrow[y_0,\dots,y_n]'\longrightarrow[y_0,\dots,y_n]
  \longrightarrow k\longrightarrow0
\]
\note{les exposants des deux $\Lambda$ sont surchargés ; on lit $n$ puis
$n+1$, ce qu'exige la suite exacte.}

Donc
\[
  (\mathcal{A}^n)_n\overset{\text{iso can}}{\simeq}(\mathcal{A}^n)_0
  \simeq k[Z_0,\dots,Z_n]
\]
(l'indice étant le deg cot), donc
\[
  \kappa_n : k[U]\longrightarrow k[Z_0,\dots,Z_n]
\]

\page{140}
\[
  s,\ s^+,\ \gamma^*,\ t
\]
\[
  (\mathcal{A}^{*}_{n})^{n}\simeq
  S\{\{X_0,\dots,X_n\}\}\big/\bigl(\textstyle\sum X_i-t\bigr)_{\mathrm{pd}}
  \otimes_S S[dX_0,\dots,dX_n]\big/\bigl(\textstyle\sum dX_i\bigr)
\]
\[
  \simeq S\{X_0,\dots,X_{n-1}\}[dX_0,\dots,dX_{n-1}]
  \xrightarrow{\ \kappa^n\ (S\text{-linéaire})\ } S
\]
\note{sous le second facteur : « en degré $n$, c'est $\simeq S$ avec la base
$dX_0\wedge\dots\wedge dX_{n-1}$ ».}

\[
  \kappa^n\bigl(X_0^{(i_0)}\cdots X_{n-1}^{(i_{n-1})}\,dX_0\wedge\dots\wedge dX_{n-1}\bigr)
  =\text{\struck{\ill{}}}\ t^{(n+\sum i_\alpha)}\,\gamma_{i_0,\dots,i_{n-1}}
\]
\note{au-dessus de l'exposant $n+\sum i_\alpha$ il écrit $\sum(i_\alpha+1)$.}

où $\gamma_{i_0\dots i_{n-1}}$ est un entier \uncertain{canonique}
\[
  \gamma_{i_0\dots i_{n-1}}=\frac{(n+\sum i_\alpha)!}{\prod_\alpha(i_\alpha!)}
  \int\!\!\cdots\!\!\int_{\sigma_n}X_0^{i_0}\cdots X_{n-1}^{i_{n-1}}\,
  dX_0\wedge\dots\wedge dX_{n-1}
\]
\[
  \sigma_n=\{X_0,\dots,X_{n-1}\geqslant0,\ X_0+\dots+X_{n-1}\leqslant1\}
\]
\note{sous l'intégrale, en accolade :
$0\leqslant\ \cdot\ \leqslant\frac{1}{n!}\gamma_{0,0,\dots,0}\overset{?}{=}\frac{1}{n!}$.}

On peut aussi calculer par récurrence sur $n$, en utilisant Stokes
\[
  \kappa^{n}\bigl(\underbrace{X_0^{(i_0)}\cdots X_{n-1}^{(i_{n-1})}\,
  dX_0\wedge\dots\wedge dX_{n-1}}_{d\left(X_{n-1}^{(i_{n-1}+1)}\,
  dX_0^{(i_0+1)}\wedge\dots\wedge dX_{n-2}^{(i_{n-2}+1)}\right)}\bigr)
  =\kappa^{n-1}\bigl(\partial_{ss}\,X_{n-1}^{(i_{n-1}+1)}\,
  dX_0^{(i_0+1)}\cdots\bigr)
\]
\note{sous le second membre, en accolade :
$X_{n-1}^{(i_{n-1}+1)}\,dX_0^{(i_0+1)}\cdots d\bigl(X_{n-2}^{(i_{n-2}+1)}\bigr)$.}

i.e. \struck{$\gamma_{i_0,\dots,i_{n-1}}$} $=$
\[
  \kappa_{n-1}\Bigl(\bigl(t-(X_0+\dots+X_{n-2})\bigr)^{(i_{n-1}+1)}
  X_0^{(i_0)}\cdots X_{n-2}^{(i_{n-2})}\,dX_0\wedge\dots\wedge dX_{n-2}\Bigr)
\]
où
\[
  \bigl(t-(X_0+\dots+X_{n-2})\bigr)^{(i_{n-1}+1)}
  =\sum_{\lambda+l_0+\dots+l_{n-2}=i_{n-1}+1}
  (-1)^{\sum l_\alpha}X_0^{(l_0)}\cdots X_{n-2}^{(l_{n-2})}\,t^{(\lambda)}
\]
\begin{align*}
  &=\kappa_{n-1}\Bigl(\sum_{\lambda+\sum_0^{n-2}l_\alpha=i_{n-1}+1}
    (-1)^{\sum l_\alpha}c_{i_0,l_0}c_{i_1,l_1}\cdots c_{i_{n-2},l_{n-2}}\\
  &\qquad\qquad t^{(\lambda)}X_0^{(i_0+l_0)}\cdots X_{n-2}^{(i_{n-2}+l_{n-2})}\,
    dX_0\wedge\dots\wedge dX_{n-2}\Bigr)\\
  &=t^{(n+\sum i_\alpha)}\Bigl[(-1)^{i_{n-1}+1}
    \sum_{\lambda+\sum l_\alpha=i_{n-1}+1}(-1)^{\sum l_\alpha}
    c_{i_0,l_0}\cdots c_{i_{n-2},l_{n-2}}\\
  &\qquad\qquad c_{\lambda,\,n-1+\sum_0^{n-2}(l_\alpha+i_\alpha)}\,
    \gamma_{i_0+l_0,\dots,i_{n-2}+l_{n-2}}\Bigr]
\end{align*}
\note{les $c_{i,l}$ sont les coefficients du produit de puissances divisées,
$X^{(i)}X^{(l)}=c_{i,l}X^{(i+l)}$, ce que la page ne dit pas ; le facteur
$c_{\lambda,\dots}$ est ajouté en accolade au-dessus de $t^{(\lambda)}$, avec
l'indice $\sum_0^{n-2}i_\alpha+l_\alpha+1$. Plusieurs exposants sont
surchargés ou biffés : $(-1)^{i_{n-1}+1}$ remplace un premier essai.}

Donc
\[
  \begin{aligned}
  \gamma_{i_0,\dots,i_{n-1}}&=\text{\struck{$(-1)^{i_{n-1}+1}$}}
  \sum_{\lambda+\sum_0^{n-2}l_\alpha=i_{n-1}+1}
  (-1)^{\sum_0^{n-2}l_\alpha}c_{i_0,l_0}\cdots c_{i_{n-2},l_{n-2}}\\
  &\qquad c_{\lambda,\,n-1+\sum_0^{n-2}(l_\alpha+i_\alpha)}\,
  \gamma_{i_0+l_0,\,i_1+l_1,\dots,\,i_{n-2}+l_{n-2}}
  \end{aligned}
\]
\note{la formule de récurrence s'arrête au bas de la page ; la suite, si elle
existe, est hors du lot.}
\end{document}
